\section*{Аннотация}

Работа посвящена оценке прогностической полезности блоков открытых характеристик компаний Y Combinator. Исследуемую совокупность составили 3\,110 компаний, активных 13 июля 2023 года. Их характеристики сопоставлены со статусами в августовской версии каталога 2026 года, скачанной автором в середине августа. Для 3\,046 компаний последующий статус известен; 272 из них поглощены или имеют публичный статус.

Описаны происхождение и ограничения данных, выполнены предобработка и разведочный анализ. Постоянный прогноз, регуляризованная логистическая регрессия и CatBoost сопоставлены во вложенной перекрёстной проверке с пятью групповыми и пятью случайными повторами. Оценены добавление и исключение блоков, ранжирование редкого исхода, веса классов при одинаковом критерии настройки и чувствительность к исключению IK12 и публичных компаний. Для парных изменений ошибки приведены условные диапазоны повторных выборок компаний и целых наборов.

Полная логистическая модель имеет LogLoss \resFullLoss{}, что на \resRelativeGain{}\% меньше потери постоянного прогноза. ROC-AUC составляет \resFullAUC{}, AP --- \resFullAP{} при частоте положительного класса 0,0893. Добавление тематических меток к модели с набором, возрастом и командой уменьшает LogLoss на \resTagGain{}. Результат указывает на ограниченный прогностический сигнал и остаётся разведочным. Условные диапазоны не учитывают переобучение; независимая внешняя проверка и оценка практической полезности не выполнены. Выводы не интерпретируются как причинные эффекты или готовые рекомендации по отбору компаний.

\addcontentsline{toc}{section}{Аннотация}
\section*{Ключевые слова}
Стартапы, Y Combinator, прогнозирование, бинарная классификация, логистическая регрессия, CatBoost, тематические метки, вложенная перекрёстная проверка, разведочный анализ данных
\pagebreak
